% ============================================================
% Sección 2: Marco Teórico
% ============================================================

\section{Marco Teórico}
\label{sec:teoria}

\subsection{Ecuación de Helmholtz para el campo óptico}
\label{subsec:helmholtz}

La propagación de luz coherente en un medio homogéneo e isótropo está gobernada
por la ecuación de Helmholtz escalar:
%
\begin{equation}
  \nabla^2 E(\mathbf{r}) + k^2 E(\mathbf{r}) = 0,
  \label{eq:helmholtz}
\end{equation}
%
donde $E(\mathbf{r}) \in \mathbb{C}$ es la amplitud compleja del campo eléctrico,
$k = 2\pi/\lambda$ es el número de onda y $\lambda$ la longitud de onda del láser.
En dos dimensiones espaciales, con $\mathbf{r} = (x, y)$, la ecuación toma la forma:
%
\begin{equation}
  \frac{\partial^2 E}{\partial x^2} + \frac{\partial^2 E}{\partial y^2} + k^2 E = 0.
  \label{eq:helmholtz2d}
\end{equation}

Como las redes neuronales operan sobre números reales, el campo complejo se descompone
en sus partes real e imaginaria: $E = E_{\mathrm{real}} + i\,E_{\mathrm{imag}}$.
Ambas componentes satisfacen la \cref{eq:helmholtz2d} de forma independiente:
%
\begin{equation}
  \nabla^2 E_{\mathrm{real}} + k^2 E_{\mathrm{real}} = 0, \qquad
  \nabla^2 E_{\mathrm{imag}} + k^2 E_{\mathrm{imag}} = 0.
  \label{eq:helmholtz_split}
\end{equation}
%
La intensidad del speckle, observable físicamente, se calcula como:
%
\begin{equation}
  I(x,y) = |E|^2 = E_{\mathrm{real}}^2 + E_{\mathrm{imag}}^2.
  \label{eq:intensidad}
\end{equation}

\subsection{Redes Neuronales Informadas por Física (PINNs)}
\label{subsec:pinns}

Las PINNs, introducidas por \citet{Raissi2019_PINNs} y revisadas extensamente en
\citet{Karniadakis2021_review} y \citet{Cuomo2022_review}, aproximan la solución de una
EDP mediante una red neuronal $\mathcal{N}_\theta(\mathbf{r})$ cuyos parámetros $\theta$
se optimizan minimizando una función de pérdida compuesta:
%
\begin{equation}
  \mathcal{L}(\theta) = \mathcal{L}_{\mathrm{datos}}(\theta)
                      + \lambda\,\mathcal{L}_{\mathrm{física}}(\theta),
  \label{eq:loss_total}
\end{equation}
%
donde $\mathcal{L}_{\mathrm{datos}}$ penaliza el error en las condiciones de frontera y
$\mathcal{L}_{\mathrm{física}}$ penaliza la violación de la EDP en el interior del dominio.
Para la ecuación de Helmholtz 2D con campo complejo:
%
\begin{align}
  \mathcal{L}_{\mathrm{datos}}(\theta) &=
    \frac{1}{N_b}\sum_{j=1}^{N_b}
    \left\|\mathcal{N}_\theta(\mathbf{r}_j^b) - E(\mathbf{r}_j^b)\right\|^2,
  \label{eq:loss_datos} \\
  \mathcal{L}_{\mathrm{física}}(\theta) &=
    \frac{1}{N_c}\sum_{i=1}^{N_c}
    \left[R_{\mathrm{real}}(\mathbf{r}_i^c)^2 + R_{\mathrm{imag}}(\mathbf{r}_i^c)^2\right],
  \label{eq:loss_fisica}
\end{align}
%
donde $R_{\mathrm{real}} = \nabla^2 E_{\mathrm{real}} + k^2 E_{\mathrm{real}}$ y
$R_{\mathrm{imag}} = \nabla^2 E_{\mathrm{imag}} + k^2 E_{\mathrm{imag}}$ son los
residuos de Helmholtz, $\{\mathbf{r}_j^b\}_{j=1}^{N_b}$ son los puntos de frontera
y $\{\mathbf{r}_i^c\}_{i=1}^{N_c}$ son los puntos de colocación interiores.
Las derivadas $\nabla^2 E$ se obtienen mediante diferenciación automática \citep{Baydin2018_autodiff}.

\subsection{Redes SIREN}
\label{subsec:siren}

Las activaciones estándar (ReLU, tanh) producen derivadas de segundo orden inestables y
exhiben \textit{spectral bias}: tienden a aprender frecuencias bajas antes que las altas,
lo que dificulta capturar los detalles finos del patrón de speckle.
\citet{Sitzmann2020_SIREN} proponen las redes SIREN (\textit{Sinusoidal Representation
Networks}), que utilizan la activación $\phi(z) = \sin(\omega_0 z)$ con una inicialización
específica que preserva la distribución de activaciones a través de todas las capas.

La arquitectura SIREN con $L$ capas ocultas de dimensión $d$ se define como:
%
\begin{equation}
  \mathbf{h}_0 = \mathbf{x}, \qquad
  \mathbf{h}_\ell = \sin\!\left(\omega_0\,(\mathbf{W}_\ell\,\mathbf{h}_{\ell-1} + \mathbf{b}_\ell)\right),
  \quad \ell = 1, \ldots, L,
  \label{eq:siren}
\end{equation}
%
donde $\omega_0$ es la frecuencia de la capa de entrada y las matrices $\mathbf{W}_\ell$
se inicializan según $\mathcal{U}(-\sqrt{6/d},\,\sqrt{6/d})$ para capas intermedias
y $\mathcal{U}(-1/n_{\mathrm{in}},\,1/n_{\mathrm{in}})$ para la primera capa
\citep{Sitzmann2020_SIREN}.
Esta inicialización garantiza que las distribuciones de activación sean aproximadamente
uniformes en $[-1, 1]$, condición necesaria para la estabilidad de derivadas de orden
arbitrario.

\subsection{Estadística del speckle completamente desarrollado}
\label{subsec:speckle}

Cuando una superficie rugosa con variación de fase $\Delta\phi \gg 2\pi$ es iluminada
por luz coherente, las contribuciones de los distintos puntos de la superficie interfieren
constructiva y destructivamente de forma aleatoria.
Por el teorema del límite central, el campo complejo resultante
$E = E_{\mathrm{real}} + i\,E_{\mathrm{imag}}$ sigue una distribución gaussiana circular
compleja \citep{Goodman2007_Speckle}:
%
\begin{equation}
  E_{\mathrm{real}},\; E_{\mathrm{imag}} \sim \mathcal{N}(0, \sigma_E^2),
  \quad \text{independientes.}
  \label{eq:campo_gaussiano}
\end{equation}

La intensidad $I = |E|^2$ sigue entonces una distribución exponencial negativa:
%
\begin{equation}
  p(I) = \frac{1}{\langle I \rangle}\exp\!\left(-\frac{I}{\langle I \rangle}\right),
  \quad I \geq 0,
  \label{eq:pdf_intensidad}
\end{equation}
%
con valor esperado $\langle I \rangle = 2\sigma_E^2$ y desviación estándar
$\sigma_I = \langle I \rangle$.
El \textit{contraste del speckle} es:
%
\begin{equation}
  C = \frac{\sigma_I}{\langle I \rangle} = 1,
  \label{eq:contraste}
\end{equation}
%
que vale exactamente 1 para speckle completamente desarrollado y disminuye
cuando el campo pierde coherencia o cuando la rugosidad es insuficiente.
Este criterio estadístico es la métrica de validación principal en NB03,
complementada con el test de Kolmogorov-Smirnov (KS) contra la distribución
exponencial teórica.
